\section{Metodologia}\label{sec:rn_metodo}

\textbf{Protocolo.} Os dados são divididos em treino (80\%, \rnTreino{} pedidos) e teste (20\%, \rnTeste), mantendo a proporção de aprovados, e o procedimento é repetido com 10 seeds (0 a 9), que mudam a divisão, as partes da validação cruzada e a semente dos modelos. Como a base não tem linhas repetidas, a divisão é estratificada simples. Quando um modelo tem hiperparâmetros a escolher, uma busca em grade com validação cruzada estratificada de 5 partes, feita só dentro do treino, escolhe a combinação de maior acurácia. No teste, reportam-se acurácia, F1 macro, a área sob a curva ROC (AUC) e, para a rede e a árvore, a matriz de confusão. O desbalanceamento é tratado pela divisão e pela validação estratificadas e por métricas que não favorecem a classe maior: o F1 macro, média simples do F1 das duas classes, a AUC, que não depende da proporção entre elas, e a matriz de confusão, que mostra o erro em cada uma. Não usamos reamostragem nem pesos por classe, porque a classe menor tem \rnRejeitadosBase{} pedidos e os modelos acertam bem as duas classes. A AUC usa a probabilidade prevista, e não a classe: é a chance de um pedido aprovado, sorteado ao acaso, receber probabilidade maior que um rejeitado. Ela mede, portanto, se o modelo ordena bem os pedidos, independentemente do corte em 50\%.

\textbf{Modelos de referência.} Cinco modelos servem de régua: a classe majoritária; a \emph{regra de um corte}, uma árvore de profundidade 1 que escolhe sozinha, no treino, o melhor atributo e o melhor limiar (escolheu o score e 549,5 nas 10 seeds); a regressão logística com padronização, que equivale a uma rede de um único neurônio; a árvore de decisão com busca em grade (critério, profundidade máxima e mínimo por folha, \rnCombinacoesArv{} combinações); e a Random Forest com 300 árvores.

\textbf{A rede neural.} A rede é um MLP do scikit-learn~\cite{pedregosa:11}: camadas ocultas com ativação ReLU~\cite{nair:10}, um neurônio de saída com sigmoide, que dá a probabilidade de aprovação, perda de entropia cruzada e otimizador Adam~\cite{kingma:15}. Antes da rede, cada atributo é padronizado (subtrai-se a média e divide-se pelo desvio padrão do treino, nunca do teste). O treino usa parada antecipada~\cite{prechelt:98}: 10\% do treino ficam de fora como validação, e o treino para quando a acurácia nessa validação passa um número fixo de épocas, a \emph{paciência}, sem melhorar. Usamos paciência de 30 épocas em vez das 10 do padrão do scikit-learn, pelo motivo mostrado na Seção~\ref{sec:rn_resultados}. A busca em grade cobre a arquitetura (8, 16, 32 ou 64 neurônios numa camada; 32 e 16 ou 64 e 32 em duas), a regularização L2 ($\alpha$ de $10^{-4}$ ou $10^{-2}$) e a taxa de aprendizado ($10^{-3}$ ou $10^{-2}$), \rnCombinacoes{} combinações.

\textbf{Tempo de execução.} Para comparar algoritmos, e não o número de núcleos que cada um usa, tudo roda num único núcleo (\texttt{n\_jobs=1} e as bibliotecas de álgebra linear limitadas a uma \emph{thread}). Antes da primeira medição, uma rodada de aquecimento é feita e descartada. Nos modelos com busca em grade, o tempo da busca é medido à parte, e o modelo final é treinado de novo com os hiperparâmetros escolhidos para medir só o treino. O resumo usa a mediana das 10 seeds. A máquina é um AMD Ryzen 5 9600X (6 núcleos) com 31 GB de memória, Windows 11, Python 3.14 e scikit-learn 1.9.1, sem outras tarefas pesadas durante a execução.
